%% ---------------------------------------------------------------------------
%% Revised Section 3 (complete). The preamble and 3.1 are condensed for the
%% 35-page limit; 3.2 and the first paragraph of 3.3 are unchanged; the
%% LLM-as-teacher paragraph and 3.4-3.6 are the revised versions.
%% ---------------------------------------------------------------------------
\section{Methodology} \label{methodology}

The framework (Figure~\ref{Figure_2}) treats text classification as graph classification so that graph and token explanations can be compared on the same decision function. Each instance is converted into four graph representations (Section~\ref{data_preprocessing}); a fine-tuned LLM supplies node features, predictions and supervisory labels (Section~\ref{fine-tuning}) for a GNN surrogate trained under the LLM-as-teacher paradigm (Section~\ref{gnn}); Shapley-based explainers are applied to the surrogates and to the LLM under a common budget (Section~\ref{explainability}); and four families of explanation metrics (Section~\ref{sec:evaluation_framework}) feed the error detectors and controls of Section~\ref{logistic_regression}. Figure~\ref{fig:explainability_pedagogic} shows how the explainability modules connect to the detectors. The modules are interchangeable, so the study can be repeated with other language models, GNN architectures, graph constructions or explainers. The code for reproducing all experiments is available via GitHub\footnote{\href{https://anonymous.4open.science/r/Graph-Neural-Networks-Enable-Superior-Error-Detection-in-NLP-Explainability-than-Language-Models-A792/README.md}{Repository for Experiment Reproducibility}}

\begin{figure}[htbp]
    \centering
    % Architecture figure on top
    \includegraphics[width=\linewidth]{Figure_1.pdf}
    \caption{End-to-end pipeline: text instances are converted to multiple graph types, embedded via fine-tuned LLM, trained as GNN surrogates via Graph Convolutional layers, and evaluated through Shapley-based explainability under unified fairness constraints.}
    \label{Figure_2}

    % Small vertical space between figures
    \vspace{1em}

    % New pedagogical explainability figure below
    \includegraphics[width=\linewidth]{Figure_2.pdf}
    \caption{Each trained model is wrapped with their respective explainability module. From these modules, we obtain agnostic metrics that can be compared across modules, those metrics are used for train each logistic regression model for error detection based on our base model prediction.}
    \label{fig:explainability_pedagogic}
\end{figure}

\subsection{Text-to-Graph Conversion} \label{data_preprocessing}

Each sentence is transformed into two hierarchical and two non-hierarchical graphs. \textbf{Hierarchical graphs} are the constituency and dependency trees produced by Stanza~\cite{qi2020stanza}: nodes are words enriched with POS tags, lemmas and morphological features (plus phrase nodes in constituency trees) and edges encode phrase-structure or head--dependent relations. \textbf{Non-hierarchical graphs} capture distributional context: window graphs connect tokens within a fixed context window and skip-gram graphs connect tokens separated by fixed intervals; nodes carry POS annotations and edges represent proximity or co-occurrence. All graphs use unique node identifiers and topology-appropriate directionality (Figure~\ref{Figure_2}, Section~1); the parser models are listed in the Supplementary Material.

\subsection{Language Model Fine-tuning and Embedding Generation} \label{fine-tuning}

Task-adapted embeddings are extracted from Transformer models fine-tuned on classification targets. We use models with [CLS] tokens to capture global sequence semantics, which also serve as embeddings for special tokens in constituency graphs, ensuring consistent handling of structural information. Models such as BERT~\cite{devlin2019bert} effectively encode both sentence-level and word-level features (see Figure~\ref{Figure_2}, Section 2).

\subsection{Graph Neural Network Training} \label{gnn}

Node features come from the fine-tuned language model. Each word node receives the mean of the last-layer hidden states of the subword tokens that make up the word; each phrase node of a constituency tree receives the hidden state of the [CLS] token. The same features are used for all graph types.

\textbf{Training under the LLM-as-teacher paradigm.} The GNN is trained to replicate the fine-tuned LLM's predictions using Graph Convolutional Network (GCN) layers~\cite{kipf2017semisupervisedclassificationgraphconvolutional}, with the LLM's classifications as supervisory labels and its hidden states as node features, following graph-aware distillation on textual graphs~\cite{mavromatis2023traingnnteachergraphaware} and LLM-to-GNN knowledge transfer~\cite{li2024enhancinggraphneuralnetworks,hu2024largelanguagemodelmeets}. Labels and node features come from the same fine-tuned checkpoint. Training details are given in Section~\ref{configuration}.

The teacher--student setup serves as an experimental control: it fixes the decision function and the input semantics shared by the graph and the token models, so that differences between their explanations can be attributed to the representation or to the explainer rather than to two independently trained classifiers. The price of this control is that the surrogate is not the original model, and what it inherits from its teacher, its errors, its confidence or its attributions, is itself one of the questions of this study. We therefore evaluate error detection under two definitions of ``error'' and add controls that isolate what the surrogate contributes (Section~\ref{logistic_regression}).

\subsection{Post-Hoc Explainability}
\label{explainability}

Post-hoc explainability techniques identify which input elements have the most influence on a trained model. We use three Shapley-based explainers---SubgraphX~\cite{yuan2021explainabilitygraphneuralnetworks}, GraphSVX~\cite{duval2021graphsvxshapleyvalueexplanations} and TokenSHAP~\cite{goldshmidt2024tokenshapinterpretinglargelanguage}---which share a common attribution principle while differing in their search strategy (see Figure~\ref{Figure_2}, Section~4).

\subsubsection{Explainability Modules}

SubgraphX uses Monte Carlo tree search (MCTS) to find a connected subgraph, within a node budget, whose Shapley-style score is highest; GraphSVX attributes importance to individual nodes from sampled coalitions and selects the top-ranked nodes; TokenSHAP estimates token-level Shapley values for the language model and we aggregate them at the word level for structural alignment with the graph explainers. Both graph explainers are applied to all four graph types, so that a signal can be attributed to the structure of the graph or to the explainer itself. SubgraphX requires a connected graph, which holds for all four; GraphSVX samples coalitions of nodes and imposes no requirement on the graph structure. Together with TokenSHAP on the LLM this yields 18 explainer--representation configurations across the two datasets.

Every explainer targets the \emph{predicted} class of the explained model, so that all confidence-based metrics refer to the class the model actually chose.

\subsubsection{Cross-Model Fair Explainability} \label{sec:crossfair}

Our experimental design adapts GraphFramEx~\cite{amara2022graphframex} fairness principles to multimodal settings following the M4 Benchmark~\cite{10.5555/3666122.3666203}. We standardise five dimensions across architectures: budget (2000 forward passes per instance), sparsity (the explanation covers 20\% of the input elements: the node budget of the MCTS search, the top-\(k\) nodes of GraphSVX and the top-\(k\) words of TokenSHAP), aligned context windows (full self-attention for LLMs \(\Leftrightarrow\) two-hop receptive fields for GNNs), method-adapted sampling strategies (polynomial approximation for TokenSHAP, coalition sampling for GraphSVX, MCTS-guided expansion for SubgraphX), and atomic-level granularity. For SubgraphX the explanation is the highest-scoring coalition within the node budget, and the node ranking required by the progression metrics of Section~\ref{sec:evaluation_framework} is derived from the search. The Supplementary Material details the specifications and nomenclature for each explainer, including this ranking.

\subsection{Multidimensional Evaluation Framework} \label{sec:evaluation_framework}

The framework characterises explanations through four families of metrics computed per instance (see Figure~\ref{Figure_2}, Section~5). All metrics use the confidence of the explained model in its predicted class \(\hat{y}\), denoted \(c(\cdot)\), under three inputs: the original input \(x\), the input restricted to the explanation \(x_{S}\) (all other elements zero-filled) and the input with the explanation removed \(x_{\setminus S}\).

\subsubsection{Stratified Evaluation Protocol}

Every metric is computed for every combination of (1) explainer (SubgraphX, GraphSVX, TokenSHAP), (2) input representation (constituency, syntactic, window, skip-gram, tokens), (3) dataset (AG News, SST-2) and (4) prediction correctness. Stratification by class is used for reporting and, in the error detectors of Section~\ref{logistic_regression}, only by the \emph{predicted} class, which is available at inference time. Core terminology is standardised in the Supplementary Material (Table~S2).

\subsubsection{Evaluated Dimensions}

Table~\ref{tab:evaluation_dimensions} summarises the four dimensions. Dimensions~1 and~2 describe confidence trajectories under progressive perturbation, Dimension~1 through the area under the curve (AUC) of the confidence as elements are deleted or inserted in decreasing order of importance, following the deletion and insertion metrics of Petsiuk et al.~\cite{petsiuk2018rise} restricted to the explanation budget and Dimension~2 through the drops at fixed numbers of elements; Dimensions~3 and~4 describe the sufficiency and necessity of the explanation through margins and confidence, respectively. Formal definitions are given in the Supplementary Material (Table~S3).

\begin{table}[!tbp]
\centering
\scriptsize
\setlength{\tabcolsep}{4pt}
\renewcommand{\arraystretch}{1.1}
\caption{Four-dimensional evaluation framework for cross-architecture explainability comparison. \(c(\cdot)\) is the confidence of the explained model in its predicted class; \(m(\cdot)\) the margin between the probability of the predicted class and the highest probability among the other classes on the same input.}
\label{tab:evaluation_dimensions}
\begin{tabular}{@{}p{0.20\textwidth}p{0.44\textwidth}p{0.30\textwidth}@{}}
\toprule
\textbf{Dimension} & \textbf{Metrics} & \textbf{Reading} \\
\midrule
\textbf{1. AUC discriminative capacity}
& Deletion AUC and insertion AUC of \(c(\cdot)\) over the importance ranking; error and correctness detection rates at thresholds \([0.1,\ldots,1.0]\)
& Low deletion AUC: the ranking removes what matters first; high insertion AUC: the top-ranked elements suffice early \\
\midrule
\textbf{2. Feature ranking stability}
& Mask-out drop \(c(x)-c(x_{\setminus S_k})\) and sufficiency drop \(c(x)-c(x_{S_k})\) for \(k\in\{1,3,5,10\}\); concentration of the mask-out drop at \(k\)
& Steep curves: concentrated importance; flat curves: distributed importance \\
\midrule
\textbf{3. Consistency across outcomes}
& Margin of the explanation alone, \(m(x_S)\), and margin of the input with the explanation removed, \(m(x_{\setminus S})\), next to the original margin \(m(x)\)
& \(m(x_S)>0\): the explanation alone reproduces the decision; \(m(x_{\setminus S})<0\): removing the explanation flips it \\
\midrule
\textbf{4. Behavioural faithfulness}
& Sufficiency \(F^{+}=c(x_S)/c(x)\), necessity \(F^{-}=1-c(x_{\setminus S})/c(x)\), asymmetry \(A=(F^{-}-F^{+})/(F^{-}+F^{+})\)
& Faithful explanations are both sufficient and necessary; \(A>0\) when the explanation is more necessary than sufficient, \(A<0\) in the opposite case \\
\bottomrule
\end{tabular}
\end{table}

\subsection{Error Signal Analysis} \label{logistic_regression}

The error-detection question is whether the metrics of Section~\ref{sec:evaluation_framework}, computed from the explanation of a prediction, allow that prediction to be classified as correct or incorrect \emph{without access to the label}. We treat it as a binary classification problem and evaluate it against the baselines that any error detector must beat.

\paragraph{Targets}
An error can be defined against two references, and we report both because they are different tasks. A \emph{gold error} is a prediction that differs from the dataset label; this definition applies to every model. A \emph{teacher disagreement}, defined for the GNN surrogates only, is a prediction that differs from the label of the teacher the surrogate was trained on, which is the target the surrogates were optimised for.

\paragraph{Feature sets}
For each instance \(i\) we build a standardised feature vector from the four dimensions: deletion and insertion AUC (Dimension~1), mask-out and sufficiency drops at \(k\in\{1,3,5,10\}\) and their concentrations (Dimension~2), the margins of the explanation alone and of its complement (Dimension~3) and \(F^{+}\), \(F^{-}\) and \(A\) (Dimension~4). The confidence-based baseline uses three label-free statistics of the explained model's output distribution: the maximum softmax probability (MSP), the margin \(m(x)\) of Table~\ref{tab:evaluation_dimensions} and the predictive entropy. Let \(\mathbf{c}_i\) denote the standardised confidence features of instance \(i\), \(\mathbf{e}_i\) its standardised explanation features and \(y_i=1\) if the prediction is correct. The detector is the logistic model
\begin{equation}
P(y_i=1\mid \mathbf{c}_i,\mathbf{e}_i)=\sigma\!\left(\beta_0+\boldsymbol{\beta}_c^{\top}\mathbf{c}_i+\boldsymbol{\beta}_e^{\top}\mathbf{e}_i\right),
\label{eq:detector}
\end{equation}
fitted by maximum likelihood with each instance weighted inversely to the frequency of its class. Three nested detectors are compared: confidence only (\(\boldsymbol{\beta}_e=0\)), explanation only (\(\boldsymbol{\beta}_c=0\)) and confidence plus explanation. The quantity of interest is the gain of the full model over the confidence-only model,
\begin{equation}
\Delta\mathrm{AUROC}=\mathrm{AUROC}(\mathbf{c},\mathbf{e})-\mathrm{AUROC}(\mathbf{c}),
\label{eq:delta}
\end{equation}
computed on out-of-fold predictions: it measures the information the explanation carries about correctness beyond what the model's confidence already provides.

\paragraph{Detectors and evaluation}
The detector of Equation~\ref{eq:detector} is trained per combination of dataset, explainer and representation, both as a pooled model and as one model per \emph{predicted} class, and evaluated by 10-fold stratified cross-validation. We report balanced accuracy, the area under the ROC curve (AUROC), the area under the precision--recall curve of the error class (AUPRC) and the Matthews correlation coefficient (MCC), together with accuracy and the majority-class baseline. Confidence intervals for \(\Delta\mathrm{AUROC}\) (Equation~\ref{eq:delta}) come from a paired bootstrap over the out-of-fold predictions: the test instances are resampled with replacement 1000 times, each instance carrying the predictions of both detectors, and the 2.5th and 97.5th percentiles of the resampled differences give the 95\% interval. The same procedure gives every interval for a difference between two models in this paper. Coefficient intervals come from 200 bootstrap refits. A gradient-boosting detector with the same features checks that the conclusions do not depend on linearity. Calibration is measured by the expected calibration error (ECE)~\cite{guo2017calibration}: the predictions are grouped into ten equal-width confidence bins, and ECE is the average over bins, weighted by their size, of the absolute difference between the mean confidence and the accuracy observed in the bin.

\paragraph{Controls for the surrogate}
Four additional analyses isolate what the surrogate contributes. (i) \emph{Cross-model detection}: a detector of BERT's gold errors built on BERT's confidence is extended, under the same nested protocol, with the surrogate's confidence, its disagreement with BERT, its explanation features, and committees of the four topologies. (ii) \emph{Explanation overlap}: for every instance, the words ranked highest by the graph explainer and by TokenSHAP at the 20\% budget are compared by Jaccard index and by the Spearman correlation of the shared words, and the overlap is added as a feature to the nested detectors. (iii) \emph{Gold-trained control}: every surrogate is retrained on the gold labels, and once more on the teacher labels with a different seed, with identical graphs, architecture and training protocol; for each model we report the fraction of BERT's errors it reproduces, the Jaccard index of the two error sets and the fraction of BERT's errors on which it is confident (\(>0.9\)) and agrees with BERT, with paired-bootstrap intervals for the difference between the gold-trained and the distilled model. (iv) \emph{Node-feature ablation}: the whole pipeline (surrogates, gold-trained controls, explainers and detectors) is repeated with node features from the pre-trained encoder instead of the fine-tuned teacher.
